\chapter{\babLima}
\label{bab:simpulan}
% Keputusan H1--H4 merujuk results/scaling_eval.json, pairing_eval.json,
% forecast_eval.json, restock_eval.json, dan results/external/E1--E5*.json.
% Keputusan H5 merujuk results/assistant_eval.json (penilai versi 2.1).

\section{Simpulan}
\label{sec:simpulan}

Penelitian ini menghasilkan Takar, aplikasi web yang menghubungkan resep standar, SOP, perhitungan variasi pesanan, panduan seduh beranimasi, pencatatan produksi, dan buku besar persediaan. Resep menyimpan ukuran, komponen, aturan takaran, parameter mutu, dan langkah seduh; produksi yang dicatat diuraikan menjadi pemakaian bahan berdasarkan BOM. Hak pengguna dibedakan antara manajer, kepala barista, dan barista. Sebanyak 706 uji backend dan 29 kasus antarmuka otomatis lulus. Hasil tersebut menjawab kebutuhan integrasi pada rumusan masalah pertama dari sisi implementasi, tetapi belum mengukur perubahan konsistensi rasa atau lama pelatihan staf di kedai nyata.

Untuk rumusan masalah kedua, penskalaan berbasis batasan memenuhi kriteria H1 pada 525 variasi pesanan yang diuji: tidak ada \textit{shot} pecahan, luapan cangkir lebih dari 5\%, atau penyimpangan rasio lebih dari 0,2. Pembanding linier mentah menghasilkan 240 pesanan dengan \textit{shot} pecahan dan 267 luapan; pembanding linier yang dibulatkan tetap menghasilkan 282 luapan. Sebagian hasil nol melekat pada bentuk algoritma, sehingga ditandai \textit{by construction}. Sebanyak 144 pesanan masih mendapat peringatan kurang isi, yang menuntut keputusan operasional lebih lanjut. Validasi menu Starbucks menunjukkan aturan \textit{shot} utuh berguna, tetapi tabel pengecualian ukuran dibutuhkan untuk merepresentasikan pola kedai tertentu; kecocokan sempurna ketika tabel diisi dari menu yang sama bukan bukti kemampuan prediksi.

Untuk rekomendasi pasangan menu, H2 didukung pada data pengembangan IBM yang bersifat fiktif, tetapi tidak didukung pada transaksi nyata Bread Basket: selisih HR@5 FP-Growth terhadap popularitas sebesar 0,0060 mempunyai IK 95\% [-0,0028; 0,0147]. Pembanding ko-okurensi juga lebih tinggi daripada FP-Growth pada HR@5 data nyata tersebut. Hasil ini menunjukkan bahwa saringan \textit{lift} aplikasi perlu diuji lagi pada transaksi kedai tujuan sebelum dipakai sebagai dasar penawaran menu.

Pada peramalan, model \textit{gradient boosting} global mempunyai WAPE terendah pada data pengembangan Maven, tetapi H3 tidak didukung pada ihelon. WAPE per minumannya sedikit lebih rendah daripada \textit{seasonal naive} tanpa selisih yang meyakinkan, sedangkan WAPE total harian justru lebih tinggi. H4 didukung untuk kebijakan pengisian ulang berbasis prakiraan dibandingkan ROP statis pada simulasi ihelon, dengan \textit{fill rate} rata-rata 0,9605 berbanding 0,7900. Analisis kontrol menunjukkan bahwa struktur tingkat sasaran dinamis berperan penting dan model global bukan peramal terbaik dalam replay itu. Oleh karena itu, hasil simulasi tidak boleh ditafsirkan sebagai bukti bahwa satu model AI selalu memberi stok terbaik.

Sistem pakar mereproduksi perhitungan EY pada seduhan dengan massa terukur dan memperlihatkan kesesuaian arah kekuatan terhadap sebagian penilaian konsumen. Akan tetapi, kesukaan rata-rata di dalam kotak \textit{Golden Cup} tidak lebih tinggi secara meyakinkan daripada di luar kotak. Diagnosis harus dibaca sebagai bantuan yang menjelaskan alasan teknis, bukan penilaian tunggal atas preferensi pelanggan.

H5 didukung secara deskriptif pada tolok ukur internal: akurasi isi asisten E4B dengan RAG dan alat 95,0\% berbanding 15,0\% tanpa konteks, sedangkan tingkat angka tanpa dasar 0\% berbanding 85,98\% menurut penilai versi 2.1. Pada kondisi lengkap, 24 jawaban berasal dari router dan alat deterministik, 10 penolakan luar cakupan dari gerbang aturan, dan 26 jawaban dari model, sehingga hasil keseluruhan tidak menyatakan akurasi Gemma 4 sendiri. Tiga jawaban E4B yang salah seluruhnya terkait prosedur/SOP. Karena butir disusun dengan bantuan AI, belum ditinjau barista praktik, dan telah dipakai untuk memperbaiki router sebelum uji akhir, angka 95,0\% merupakan hasil regresi internal, bukan estimasi generalisasi pada kedai baru. Dengan batas itu, rumusan masalah ketiga terjawab sebagai perbandingan terukur antarmodul dan kondisi, bukan pembuktian manfaat operasional pada pengguna nyata.

Kelayakan fungsional awal didukung oleh pengujian otomatis dan pengukuran kinerja lokal, tetapi rumusan masalah keempat belum sepenuhnya terjawab. Uji penerimaan pengguna bersama barista atau pengelola dan skor SUS belum tersedia. Kesimpulan tentang kemudahan pakai, pengurangan waktu pelatihan, pengurangan limbah, maupun keuntungan usaha memerlukan pengamatan pada kedai mitra, bukan hanya hasil dataset publik dan simulasi.

\section{Saran}
\label{sec:saran}

Penerapan berikutnya sebaiknya dimulai dengan mengukur resep dan ukuran cangkir di kedai mitra, memvalidasi takaran pengisi dan ambang kurang isi bersama kepala barista, lalu mencatat produksi serta stok aktual beberapa minggu. Parameter \textit{lead time}, kemasan, kehilangan bahan, dan stok awal perlu diambil dari pemasok dan buku stok nyata. Dengan demikian, saran pemesanan ulang dapat dihitung dari kondisi yang benar-benar dihadapi kedai, bukan dari data contoh.

Model pasangan menu dan peramalan perlu dilatih serta dinilai ulang pada transaksi kedai tujuan memakai pemisahan waktu. Ko-okurensi dan model musiman sederhana harus tetap menjadi pembanding karena keduanya unggul pada sebagian validasi eksternal. Kebijakan \textit{restock} perlu dievaluasi bersama biaya pemesanan, ruang penyimpanan, masa simpan susu dan bahan lain, serta barang yang sedang dipesan. Pengembangan lanjutan dapat menambahkan pencatatan pesanan pembelian dan integrasi kasir setelah alur satu outlet stabil.

Asisten lokal perlu diaudit dengan himpunan pertanyaan baru yang disusun dan diperiksa barista independen, khususnya angka takaran, saran koreksi, pertanyaan di luar cakupan, dan sumber jawaban. Himpunan itu harus dipisahkan dari butir yang telah memandu penyempurnaan router alat agar ketepatan pada kasus baru dapat diukur. Ketepatan alat, jawaban model, hasil pencarian cadangan, dan penolakan oleh aturan perlu dilaporkan terpisah. Hasil model dan aturan pakar harus ditinjau saat resep atau SOP berubah. Akhirnya, uji tugas barista dan SUS perlu dilaksanakan dengan peserta, perangkat, dan kondisi kerja yang dicatat jelas; hasilnya kemudian dipakai untuk memperbaiki navigasi, istilah, serta alur produksi sebelum klaim manfaat operasional dibuat.
